% Algebraic Trust Composition: A Lean Mechanization of Sandbox ⊗ Cost-Guard
% Invariant Preservation for Verified MCP Tool Execution
%
% Target: PLDI 2028 (ACM SIGPLAN Conference on Programming Language Design
%         and Implementation)
% Format: ACM sigconf (acmart)
% Author: Joseph Krzywoszyja
% Draft:  June 2026
%
% Paper 29 of the HoloScript research program.
% Lean mechanization: research/paper-29-algebraic-trust-toolsandbox/ATC/
%   ATC.Basic          -- foundational definitions (SandboxPolicy, CostGuard, etc.)
%   ATC.CompositionLaw -- 23 theorems, 3 axioms, 0 sorry
%
% Cross-references: P3 (CRDT semiring), P4 (Sandbox), P8 (Unified Transform),
%                  P25 (Multi-Brain budget)
% Composition theorem: cp = sandbox ⊗ costGuard implies:
%   (1) sandbox.permits cap = false → cp.permits cap = false
%   (2) costGuard.permits = false  → cp.permits cap = false
%   (3) cp.permits cap = sandbox.permits cap ∧ costGuard.permits

\documentclass[sigconf]{acmart}

% ── Packages ──────────────────────────────────────────────────────────────────
\usepackage{amsmath}
\let\Bbbk\relax  % prevent amssymb/newtxmath clash
\usepackage{amssymb}
\usepackage{booktabs}
\usepackage{multirow}
\usepackage{xcolor}
\usepackage{listings}
\usepackage{subcaption}
\usepackage{balance}
% todonotes replacement (avoid package dependency; renders inline in draft)
\newcommand{\todo}[1]{\textbf{[TODO: #1]}}

% ── Theorem environments ──────────────────────────────────────────────────────
\newtheorem{definition}{Definition}
\newtheorem{theorem}{Theorem}
\newtheorem{lemma}{Lemma}
\newtheorem{corollary}{Corollary}
\newtheorem{property}{Property}

% ── Code listing styles ───────────────────────────────────────────────────────
\lstdefinelanguage{Lean4}{
  keywords={theorem,lemma,def,structure,inductive,axiom,namespace,end,
            import,by,intro,cases,exact,rfl,rw,simp,constructor,left,right,
            have,show,contradiction,exfalso,native_decide,unfold,match,where,
            private,with,fun,true,false,forall,exists},
  keywordstyle=\bfseries\color{blue!70!black},
  commentstyle=\itshape\color{green!50!black},
  stringstyle=\color{orange!70!black},
  morecomment=[l]{--},
  morecomment=[s]{/-}{-/},
  morestring=[b]",
  sensitive=true,
}

\lstdefinelanguage{TypeScript}{
  keywords={interface,readonly,void,null,string,number,boolean,
            Record,unknown,Promise,export,const,let,function,
            return,if,for,class,extends,implements,async,await,
            new,type,import,from,false,true},
  keywordstyle=\bfseries\color{blue!70!black},
  commentstyle=\itshape\color{green!50!black},
  stringstyle=\color{orange!70!black},
  morecomment=[l]{//},
  morestring=[b]',
  morestring=[b]",
}

\lstset{
  basicstyle=\scriptsize\ttfamily,
  breaklines=true,
  frame=single,
  numbers=left,
  numberstyle=\tiny\color{gray},
  xleftmargin=1.5em,
  aboveskip=0.5\baselineskip,
  belowskip=0.5\baselineskip,
}

% ── Notation macros ───────────────────────────────────────────────────────────
\newcommand{\compose}{\mathbin{\otimes}}
\newcommand{\permits}{\operatorname{permits}}
\newcommand{\allows}{\operatorname{allows}}
\newcommand{\overBudget}{\operatorname{isOverBudget}}
\newcommand{\mergeadd}{\mathbin{\oplus}}
\newcommand{\mergemul}{\mathbin{\otimes}}
\newcommand{\ATC}{\textsc{ATC}}
\newcommand{\SP}{\mathit{sp}}
\newcommand{\CG}{\mathit{cg}}
\newcommand{\CP}{\mathit{cp}}
\newcommand{\Bool}{\mathbb{B}}

% ── Lotus data-contract annotation ───────────────────────────────────────────
\providecommand{\measuredFrom}[1]{}

% ══════════════════════════════════════════════════════════════════════════════
\begin{document}

\title{Algebraic Trust Composition: A Lean Mechanization of Sandbox $\compose$
       Cost-Guard Invariant Preservation for Verified MCP Tool Execution}

\author{Joseph Krzywoszyja}
\email{brianonbased@gmail.com}
\affiliation{%
  \institution{Independent Researcher}
  \country{Australia}
}

% ── Abstract ──────────────────────────────────────────────────────────────────
\begin{abstract}
We present a Lean~4 mechanization of the \emph{algebraic trust composition
theorem} for AI agent tool execution.  When a \emph{sandbox policy} $\SP$ and
a \emph{cost guard} $\CG$ are combined via the composition operator
$\compose$, the resulting composed policy $\CP = \SP \compose \CG$ is
provably at least as restrictive as either component: (1)~any capability
denied by the sandbox remains denied under composition
(\emph{sandbox preservation}); (2)~any execution blocked by a budget
ceiling remains blocked under composition (\emph{budget preservation});
and (3)~permission reduces to conjunction of the two guards.  The
mechanization (\ATC{}) comprises 23 proved theorems, 3 named axioms
(externally validated by Paper~3's CRDT semiring proofs), and 0 uses
of \texttt{sorry}, verified by \texttt{lake build} under Lean~4.15.0.
Every abstract type maps to a production component of the HoloScript
MCP server; the theorem is therefore not merely a formal artifact but
a correctness certificate for deployed agent infrastructure.
\end{abstract}

\begin{CCSXML}
<ccs2012>
  <concept>
    <concept_id>10003752.10003790.10003791</concept_id>
    <concept_desc>Theory of computation~Type theory</concept_desc>
    <concept_significance>500</concept_significance>
  </concept>
  <concept>
    <concept_id>10002978.10003014</concept_id>
    <concept_desc>Security and privacy~Software and application security</concept_desc>
    <concept_significance>500</concept_significance>
  </concept>
  <concept>
    <concept_id>10003752.10003790.10003810</concept_id>
    <concept_desc>Theory of computation~Formal languages and automata theory</concept_desc>
    <concept_significance>300</concept_significance>
  </concept>
  <concept>
    <concept_id>10011007.10011006.10011008</concept_id>
    <concept_desc>Software and its engineering~General programming languages</concept_desc>
    <concept_significance>300</concept_significance>
  </concept>
</ccs2012>
\end{CCSXML}

\ccsdesc[500]{Theory of computation~Type theory}
\ccsdesc[500]{Security and privacy~Software and application security}
\ccsdesc[300]{Theory of computation~Formal languages and automata theory}
\ccsdesc[300]{Software and its engineering~General programming languages}

\keywords{type-system composition, capability restriction, sandbox policy,
  cost guard, Lean 4, mechanized proof, MCP tool execution, semiring algebra,
  AI agent security, verified enforcement}

\maketitle

% ══════════════════════════════════════════════════════════════════════════════
\section{Introduction}
\label{sec:intro}

Modern AI agents are no longer passive text processors.  They invoke tools: they
read files, call APIs, spawn processes, and spend real money on cloud resources.
The Model Context Protocol (MCP)~\cite{mcp-spec,hou2025mcp} has rapidly
become a de-facto standard for equipping language models with such capabilities.
Ensuring that an agent \emph{cannot} exceed its authorization---either by
reaching a forbidden capability class or by exhausting a budget ceiling---is
therefore a concrete engineering requirement, not a theoretical nicety.

Two complementary enforcement mechanisms address this requirement:

\begin{itemize}
  \item A \textbf{sandbox policy} ($\SP$) declares which capability classes
        (network access, file writes, compute, nested calls) an agent may use.
        Violations are denied before execution.
  \item A \textbf{cost guard} ($\CG$) tracks cumulative spend against a daily
        budget ceiling.  Once the ceiling is reached, all execution is denied
        until the daily rollover.
\end{itemize}

The central engineering question is: \emph{when these two enforcement layers
are composed, do their invariants survive?}  Specifically, could composition
somehow ``open'' a capability the sandbox forbids, or ``clear'' a budget the
cost guard has already exceeded?  It is intuitively clear that the answer
should be no, but intuition is not a proof---and enforcement code has a
history of subtle composition bugs~\cite{wahbe1993efficient,yee2009native}.

\paragraph{Contribution.}
This paper proves, mechanically in Lean~4, that the composition operator
$\compose$ \emph{is a safety intersection}: the composed policy is at least as
restrictive as either component alone.  Our contributions are:

\begin{enumerate}
  \item \textbf{Composition theorem} (\S\ref{sec:theorem}): a formal statement
        and proof that $\CP = \SP \compose \CG$ preserves both the sandbox
        invariant and the budget invariant under all inputs.

  \item \textbf{Lean mechanization} (\S\ref{sec:lean}): 23~theorems,
        3~named axioms (all externally validated), 0~\texttt{sorry},
        structured in two modules (\texttt{ATC.Basic},
        \texttt{ATC.CompositionLaw}).

  \item \textbf{Production mapping} (\S\ref{sec:mapping}): a seven-row
        correspondence table linking every abstract model element to the
        HoloScript production codebase.

  \item \textbf{Axiom audit} (\S\ref{sec:axioms}): documentation of the three
        axioms, their proof obligations, and their external validation in
        Paper~3's CRDT semiring suite~\cite{paper-3-crdt-ecoop2027}.
\end{enumerate}

\paragraph{Scope.}
This paper covers the composition theorem for Layer~1 of the W.GOLD.189
trust-by-construction stack (capability restriction + budget bounding).
Layer~2 (CAEL hash chain audit) and Layer~3 (simulation-contract oracle) are
addressed in companion papers~\cite{cael-paper,trust-by-construction} and
remain open for mechanization at the circuit level~(\S\ref{sec:conclusion}).

% ══════════════════════════════════════════════════════════════════════════════
\section{Background: HoloScript Trust Architecture}
\label{sec:background}

\subsection{MCP Tool Execution}

HoloScript exposes its simulation and compilation capabilities through
MCP~\cite{mcp-spec}, a JSON-RPC protocol in which a \emph{server} advertises
tools and a \emph{client} (the AI agent) invokes them.  Each tool call passes
through a pre-execution gate in the MCP server before any side-effectful work
begins.  This gate is the composition point that this paper reasons about.

The current production gate (in \texttt{mcp-server/src/security/}) checks, in
order:

\begin{enumerate}
  \item Does the sandbox policy permit the capability class being invoked?
  \item Does the cost guard report remaining budget?
\end{enumerate}

Both conditions must hold; either failure denies execution.  This ``both must
pass'' logic is exactly the conjunction that the composition theorem
formalizes.

\subsection{Sandbox Policies}

A sandbox policy maps each \emph{capability class} to a Boolean permission.
HoloScript recognizes four classes: \texttt{NetworkCap} (outbound connections,
rate limits, allowed hosts), \texttt{FileCap} (filesystem writes, allowed
paths, size limits), \texttt{ComputeCap} (CPU ceiling, memory ceiling,
timeout), and \texttt{NestedCap} (recursive MCP calls).  The production type
in \texttt{sandbox-policy.ts} has per-class configuration structs; for the
composition theorem we abstract these to a single Boolean \texttt{permits}
function, which captures the essential permission decision.

Sandbox composition appears in multi-brain
scenarios~\cite{paper-4-sandbox-usenix}: a fleet orchestrator may compose an
outer \emph{network sandbox} (denying all egress) with an inner \emph{compute
sandbox} (capping CPU).  The composition theorem guarantees that neither the
network denial nor the CPU cap can be weakened by composition.

\subsection{Cost Guards}

A cost guard tracks cumulative USD spend against a daily ceiling.
\texttt{CostGuard.isOverBudget()} returns \texttt{true} once
\texttt{spentUsd \(\geq\) dailyBudgetUsd}, blocking all further execution.
This is an \emph{x402} spend model~\cite{paper-3-crdt-ecoop2027}: the agent's
budget is a first-class ecosystem primitive, not a soft advisory.

For the composition theorem, the essential property is that the cost guard
either permits or denies all execution uniformly (it is capability-class-agnostic---any capability costs, no capability is exempt from the budget).

\subsection{The Simulation-as-Proof Loop}

HoloScript's design thesis is that a correct simulation \emph{is} a
proof~\cite{trust-by-construction}: a simulator that enforces its contracts at
every step provides mechanical evidence of the behaviors it produces.  Trust
composition is part of this loop: an agent that can provably never exceed its
sandbox or its budget is a weaker adversary and a stronger contributor to
verifiable simulation outcomes.  The composition theorem in this paper is a
Layer~1 instance of this thesis---the simplest building block in the
three-layer W.GOLD.189 stack.

% ══════════════════════════════════════════════════════════════════════════════
\section{Algebraic Framework}
\label{sec:algebra}

We now give the formal definitions that the Lean mechanization instantiates.
All definitions correspond directly to types in \texttt{ATC.Basic.lean}.

\begin{definition}[Capability class]
\label{def:cap}
$\mathcal{C} = \{\texttt{NetworkCap},\,\texttt{FileCap},\,\texttt{ComputeCap},\,\texttt{NestedCap}\}$
is the finite set of capability classes recognized by the MCP gate.
\end{definition}

\begin{definition}[Sandbox policy]
\label{def:sandbox}
A \emph{sandbox policy} is a function $\SP : \mathcal{C} \to \Bool$ such that
$\SP(\mathit{cap}) = \top$ means the capability is permitted and $\bot$ means
it is denied.  We write $\SP.\permits(\mathit{cap})$ for $\SP(\mathit{cap})$.

Two canonical policies are distinguished:
\begin{itemize}
  \item $\mathtt{sandboxAll}$: $\lambda\,c.\,\top$ (maximally permissive; the
        identity element for sandbox composition).
  \item $\mathtt{sandboxNone}$: $\lambda\,c.\,\bot$ (maximally restrictive;
        the annihilator).
\end{itemize}
\end{definition}

\begin{definition}[Cost guard]
\label{def:costguard}
A \emph{cost guard} is a pair $\CG = (\mathit{limit},\,\mathit{spent})$ where
$\mathit{limit} \in \mathbb{N}_\infty$ is the daily budget ceiling (with
$\infty$ meaning unlimited) and $\mathit{spent} \in \mathbb{N}$ is the
cumulative spend.  The \emph{permits} predicate is:
\[
  \CG.\permits \;=\; \begin{cases}
    \top & \text{if } \mathit{limit} = \infty \\
    \top & \text{if } \mathit{spent} < \mathit{limit} \\
    \bot & \text{if } \mathit{spent} \geq \mathit{limit}
  \end{cases}
\]
Two canonical guards are distinguished:
\begin{itemize}
  \item $\mathtt{costUnlimited}$: $\mathit{limit} = \infty,\,\mathit{spent} = 0$
        (always permits; the identity element).
  \item $\mathtt{costZero}$: $\mathit{limit} = 0,\,\mathit{spent} = 1$
        (always denies; the annihilator).
\end{itemize}
\end{definition}

\begin{definition}[Composed policy]
\label{def:composed}
Given a sandbox policy $\SP$ and a cost guard $\CG$, the \emph{composed
policy} $\CP = \SP \compose \CG$ is the structure
$\CP = (\SP,\,\CG)$ with permission predicate:
\[
  \CP.\permits(\mathit{cap}) \;=\; \SP.\permits(\mathit{cap}) \;\wedge\; \CG.\permits
\]
Execution of a capability under $\CP$ yields one of four outcomes:
\[
  \mathtt{execute}(\CP,\,c) \;=\; \begin{cases}
    \mathtt{allowed}       & \text{if } \SP.\permits(c) \wedge \CG.\permits \\
    \mathtt{deniedSandbox} & \text{if } \neg\SP.\permits(c) \wedge \CG.\permits \\
    \mathtt{deniedBudget}  & \text{if } \SP.\permits(c) \wedge \neg\CG.\permits \\
    \mathtt{deniedBoth}    & \text{if } \neg\SP.\permits(c) \wedge \neg\CG.\permits
  \end{cases}
\]
\end{definition}

\begin{definition}[Merge strategy / semiring]
\label{def:semiring}
The five \emph{merge strategies} from W.GOLD.189~\cite{paper-3-crdt-ecoop2027}
are:
\[
  \mathcal{S} = \{\texttt{minPlus},\,\texttt{maxPlus},\,\texttt{authorityWeighted},\,\texttt{domainOverride},\,\texttt{strictError}\}
\]
For each $s \in \mathcal{S}$, the semiring addition
$\mergeadd_s : \mathbb{N} \times \mathbb{N} \to \mathbb{N}$ and
multiplication $\mergemul_s : \mathbb{N} \times \mathbb{N} \to \mathbb{N}$
are:

\medskip
\begin{center}
\begin{tabular}{lll}
\toprule
Strategy & $a \mergeadd_s b$ & $a \mergemul_s b$ \\
\midrule
\texttt{minPlus}           & $\min(a,b)$      & $a + b$ \\
\texttt{maxPlus}           & $\max(a,b)$      & $a + b$ \\
\texttt{authorityWeighted} & $a + b$          & $a \cdot b$ \\
\texttt{domainOverride}    & $\max(a,b)$      & $a + b$ \\
\texttt{strictError}       & $a$              & $a$ \\
\bottomrule
\end{tabular}
\end{center}
\medskip

The min-plus and max-plus rows are tropical semirings~\cite{mohri2002semiring,green-semiring2007};
the authority-weighted row is the standard $(\mathbb{N},+,\times)$ semiring.
\end{definition}

% ══════════════════════════════════════════════════════════════════════════════
\section{Composition Theorem}
\label{sec:theorem}

\begin{theorem}[Sandbox $\compose$ cost-guard composition]
\label{thm:main}
For any composed policy $\CP = (\SP,\,\CG)$ and any capability $c \in
\mathcal{C}$, the following three properties hold simultaneously:

\begin{enumerate}
  \item \textbf{Sandbox preservation}:
        \[\SP.\permits(c) = \bot \;\Longrightarrow\; \CP.\permits(c) = \bot\]
  \item \textbf{Budget preservation}:
        \[\CG.\permits = \bot \;\Longrightarrow\; \CP.\permits(c) = \bot\]
  \item \textbf{Composed equals conjunction}:
        \[\CP.\permits(c) = \SP.\permits(c) \;\wedge\; \CG.\permits\]
\end{enumerate}
\end{theorem}

\begin{proof}[Proof sketch]
Properties~(1) and~(2) follow directly from property~(3) and the
absorption laws of Boolean conjunction:
\[
  \bot \wedge x = \bot \quad \text{and} \quad x \wedge \bot = \bot
\]
Property~(3) holds by the definition of $\CP.\permits$ (Definition~\ref{def:composed}).

In the Lean mechanization, the conjunction law~(3) is proved by \texttt{rfl}
(it follows from unfolding the definition), while properties~(1) and~(2) are
proved using \texttt{and\_false\_left} and \texttt{and\_false\_right}---private
lemmas about Boolean operations established before the main theorem.  The
conjunction of all three is stated as \texttt{composition\_preserves\_both}
(Theorem~16 in Table~\ref{tab:theorems}) and proved in four lines:
two \texttt{exact} calls reusing the preservation lemmas plus \texttt{rfl}
for property~(3).
\end{proof}

\begin{corollary}[Safety intersection]
\label{cor:safety}
The composition operator $\compose$ is a \emph{safety intersection}: the
composed policy is at least as restrictive as either component alone.  In
particular, the permission set of $\CP$ is contained in the permission set of
$\SP$ and in the permission set of the uniform-capability projection of $\CG$.
\end{corollary}

\begin{corollary}[Allowed implies both]
\label{cor:allowed}
If $\mathtt{execute}(\CP,\,c) = \mathtt{allowed}$, then $\SP.\permits(c) = \top$
and $\CG.\permits = \top$.  The contrapositive: either denial is sufficient to
block execution.
\end{corollary}

Corollary~\ref{cor:allowed} is proved as \texttt{allowed\_implies\_both\_permitted}
(Theorem~19 in the mechanization) using a four-way case split on the
execution-result discriminant; each impossible case is discharged by
\texttt{native\_decide}.

\begin{theorem}[Algebraic laws]
\label{thm:laws}
The following algebraic laws hold:

\begin{itemize}
  \item \textbf{Identity}: $\mathtt{sandboxAll} \compose \CG$ reduces to
        $\CG$'s permission; $\SP \compose \mathtt{costUnlimited}$ reduces to
        $\SP$'s permission.
  \item \textbf{Annihilator}: $\mathtt{sandboxNone} \compose \CG$ and
        $\SP \compose \mathtt{costZero}$ both deny everything.
  \item \textbf{Commutativity (sandbox intersection)}: For any two sandbox
        policies $\SP_1, \SP_2$ and any $c$,
        $\SP_1.\permits(c) \wedge \SP_2.\permits(c) = \SP_2.\permits(c) \wedge \SP_1.\permits(c)$.
  \item \textbf{Commutativity (cost conjunction)}:
        $\CG_1.\permits \wedge \CG_2.\permits = \CG_2.\permits \wedge \CG_1.\permits$.
\end{itemize}
\end{theorem}

The identity and annihilator laws correspond to Theorems~9--15 in
Table~\ref{tab:theorems}.  The commutativity laws (Theorems~7--8) follow from
Boolean commutativity by a simple two-way case split in Lean.

% ══════════════════════════════════════════════════════════════════════════════
\section{Lean Mechanization}
\label{sec:lean}

\subsection{Module Structure}

The \ATC{} project is organized in two Lean~4 modules:

\begin{description}
  \item[\texttt{ATC.Basic}]  Foundational definitions: \texttt{CapabilityClass},
    \texttt{SandboxPolicy}, \texttt{CostGuard}, \texttt{ComposedPolicy},
    \texttt{MergeStrategy}, \texttt{mergeAdd}, \texttt{mergeMul},
    \texttt{ExecutionResult}, and \texttt{execute}.  Zero theorems, zero
    axioms---definitions only.

  \item[\texttt{ATC.CompositionLaw}]  All 23~theorems and 3~axioms,
    structured in nine sections (\S1~sandbox preservation, \S2~budget
    preservation, \S3~commutativity, \S4~identity laws, \S5~annihilators,
    \S6~the main composition theorem, \S7~execution classification,
    \S8~semiring axioms, \S9~concrete strategy proofs).
\end{description}

\paragraph{Build.}

\begin{lstlisting}[language=bash,numbers=none,frame=none]
cd research/paper-29-algebraic-trust-toolsandbox
lake build   # requires leanprover/lean4:v4.15.0
\end{lstlisting}

\noindent
\texttt{lake build} exits~0 with no errors.  There are no \texttt{sorry}s to
suppress and no \texttt{native\_decide} timeouts.

\subsection{Theorem Index}

Table~\ref{tab:theorems} lists all 23~theorems and 3~axioms.  The ``Proof
technique'' column names the primary Lean tactic; many proofs combine two
short steps.

\begin{table}[t]
\caption{ATC.CompositionLaw: all 23 theorems and 3 axioms}
\label{tab:theorems}
\small
\begin{tabular}{rll}
\toprule
\# & Name & Proof technique \\
\midrule
1  & \texttt{sandbox\_preservation}                      & \texttt{Bool} def + \texttt{rfl} \\
2  & \texttt{sandbox\_preservation\_contrapositive}      & \texttt{cases} + contradiction \\
3  & \texttt{sandbox\_preservation\_for\_all\_cost\_guards} & reuse \#1 \\
4  & \texttt{budget\_preservation}                       & \texttt{Bool} def + \texttt{rfl} \\
5  & \texttt{budget\_preservation\_contrapositive}       & \texttt{cases} + contradiction \\
6  & \texttt{budget\_preservation\_for\_all\_sandbox\_policies} & reuse \#4 \\
7  & \texttt{sandbox\_composition\_commutative}          & \texttt{Bool} cases \\
8  & \texttt{cost\_composition\_commutative}             & \texttt{Bool} cases \\
9  & \texttt{sandbox\_identity\_left}                    & \texttt{cases} + \texttt{rfl} \\
10 & \texttt{cost\_identity\_permits}                    & \texttt{unfold} + \texttt{rfl} \\
11 & \texttt{cost\_identity\_composition}                & reuse \#10 + \texttt{Bool} \\
12 & \texttt{sandbox\_annihilator\_left}                 & \texttt{Bool} def + \texttt{rfl} \\
13 & \texttt{sandbox\_annihilator\_right}                & \texttt{Bool} def + \texttt{rfl} \\
14 & \texttt{cost\_annihilator\_permits}                 & \texttt{native\_decide} \\
15 & \texttt{cost\_annihilator}                          & reuse \#14 + \texttt{Bool} \\
16 & \texttt{composition\_preserves\_both}               & THE THEOREM (\#1 + \#4 + \texttt{rfl}) \\
17 & \texttt{execute\_preserves\_sandbox}                & \texttt{execute\_unfold} + \texttt{cases} \\
18 & \texttt{execute\_preserves\_budget}                 & \texttt{execute\_unfold} + \texttt{cases} \\
19 & \texttt{allowed\_implies\_both\_permitted}          & \texttt{cases} + \texttt{native\_decide} \\
20 & \texttt{semiring\_merge\_commutative}               & \emph{axiom} (P3 Theorem~1) \\
21 & \texttt{semiring\_merge\_associative}               & \emph{axiom} (P3 Theorem~2) \\
22 & \texttt{composition\_distributes\_over\_merge}      & \emph{axiom} (P3 Theorem~3) \\
23 & \texttt{min\_plus\_commutative}                     & \texttt{simp} + \texttt{Nat.min\_comm} \\
24 & \texttt{max\_plus\_commutative}                     & \texttt{simp} + \texttt{Nat.max\_comm} \\
25 & \texttt{authority\_weighted\_commutative}           & \texttt{simp} + \texttt{Nat.add\_comm} \\
26 & \texttt{min\_plus\_associative}                     & \texttt{simp} + \texttt{Nat.min\_assoc} \\
27 & \texttt{max\_plus\_associative}                     & \texttt{simp} + \texttt{Nat.max\_assoc} \\
28 & \texttt{authority\_weighted\_associative}           & \texttt{simp} + \texttt{Nat.add\_assoc} \\
\bottomrule
\end{tabular}
\end{table}

\subsection{Five Theorem Families}

The 23~theorems cluster into five families aligned with the paper's sections:

\paragraph{Family~1: Core composition (Theorems~1--6).}
These six theorems establish the composition theorem's two directions
(sandbox preservation and budget preservation) plus their contrapositives and
universally-quantified forms.  Together they prove that \emph{no capability
opens under composition} and \emph{no budget clears under composition}.

\paragraph{Family~2: Algebraic laws (Theorems~7--15).}
Commutativity of both components, plus identity and annihilator laws.
These establish that the composition operator forms a well-behaved algebraic
structure, enabling equational reasoning about policy chains.

\paragraph{Family~3: Main theorem (Theorem~16).}
\texttt{composition\_preserves\_both} is the central conjunction that
packages all three properties of Theorem~\ref{thm:main}.  It is proved in
four lines by appealing directly to Family~1 plus \texttt{rfl}.

\paragraph{Family~4: Execution classification (Theorems~17--19).}
These theorems connect the Boolean \texttt{permits} predicate to the
four-valued \texttt{ExecutionResult} discriminant, showing that sandbox
denials and budget denials are faithfully propagated into observable execution
outcomes.

\paragraph{Family~5: Semiring strategy laws (Theorems~20--28).}
Three abstract axioms (externally validated in Paper~3) plus six concrete
proofs for the min-plus, max-plus, and authority-weighted strategies.  These
establish that the merge operator is a commutative, associative semiring
addition and that composition distributes over it.

\subsection{Key Lean Excerpt}

The core of the mechanization is the main composition theorem and its two
supporting lemmas.  Below are representative extracts.

\begin{lstlisting}[language=Lean4,caption={Sandbox preservation lemma (ATC.CompositionLaw \S1)}]
theorem sandbox_preservation (cp : ComposedPolicy) (cap : CapabilityClass) :
    cp.sandbox.permits cap = false ->
    cp.permits cap = false := by
  intro h_denied
  have h : (cp.permits cap) =
           (cp.sandbox.permits cap && cp.costGuard.permits) := rfl
  rw [h, h_denied, and_false_left]
\end{lstlisting}

\begin{lstlisting}[language=Lean4,caption={Main composition theorem (ATC.CompositionLaw \S6)}]
theorem composition_preserves_both
    (cp : ComposedPolicy) (cap : CapabilityClass) :
    (cp.sandbox.permits cap = false -> cp.permits cap = false) /\
    (cp.costGuard.permits = false  -> cp.permits cap = false) /\
    (cp.permits cap = cp.sandbox.permits cap && cp.costGuard.permits) := by
  constructor
  . exact sandbox_preservation cp cap
  . constructor
    . exact budget_preservation cp cap
    . rfl
\end{lstlisting}

\begin{lstlisting}[language=Lean4,caption={Semiring axiom (externally validated)}]
axiom semiring_merge_commutative (s : MergeStrategy) (a b : Nat) :
    mergeAdd s a b = mergeAdd s b a
-- Paper 3 Theorem 1: validated on 10000/10000 pairs, commit 66d58b58
\end{lstlisting}

\noindent
The \texttt{axiom} keyword in Lean~4 introduces a hypothesis that the kernel
accepts without a proof term; the obligation is discharged externally.
\S\ref{sec:axioms} documents this discharge.

% ══════════════════════════════════════════════════════════════════════════════
\section{Production Model Mapping}
\label{sec:mapping}

The abstract model is not merely illustrative---every type maps to a
production component.  Table~\ref{tab:mapping} gives the full correspondence.

\begin{table}[t]
\caption{Abstract model to production code mapping}
\label{tab:mapping}
\small
\begin{tabular}{lll}
\toprule
Abstract element & Production construct & Source file \\
\midrule
\texttt{SandboxPolicy.permits}  & \texttt{SandboxPolicy.allows(cap)} &
  \texttt{mcp-server/src/security/}\\
  & & \quad\texttt{sandbox-policy.ts} \\[2pt]
\texttt{CostGuard.permits}      & \texttt{!CostGuard.isOverBudget()} &
  \texttt{holoscript-agent/src/}\\
  & & \quad\texttt{cost-guard.ts} \\[2pt]
\texttt{ComposedPolicy.permits} & \texttt{allows(cap) \&\& !isOverBudget()} &
  MCP server compose check \\[2pt]
\texttt{MergeStrategy}          & \texttt{strategyToSemiring()} &
  \texttt{core/src/compiler/traits/}\\
  & & \quad\texttt{Semiring.ts} \\[2pt]
\texttt{mergeAdd}               & \texttt{ProvenanceSemiring.merge()} &
  \texttt{core/src/compiler/traits/}\\
  & & \quad\texttt{ProvenanceSemiring.ts} \\[2pt]
\texttt{CAELEvent}              & CAEL pipeline events &
  \texttt{engine/src/simulation/} \\[2pt]
\texttt{ExecutionResult}        & MCP tool execution result &
  \texttt{mcp-server/src/security/}\\
  & & \quad\texttt{fork-sandbox-gate.ts} \\
\bottomrule
\end{tabular}
\end{table}

\paragraph{Abstraction gap.}
The abstract \texttt{SandboxPolicy.permits} function collapses the production
type's per-class configuration structs (allowed hosts, path lists, rate limits,
etc.) into a single Boolean decision.  The composition theorem is correct at
this level of abstraction because the conjunctive gate structure is the same
regardless of the per-class details: if the outer Boolean is \texttt{false},
execution is denied, and no inner configuration can override this.  The
finer-grained limits compose pointwise in exactly the same way.

Similarly, \texttt{CostGuard.permits} in the abstract model corresponds to the
negation of \texttt{isOverBudget()}: the abstract model exposes the permission
(true when budget remains), while the production API exposes the denial (true
when over budget).  The semantics are identical.

\paragraph{The compose check in production.}
The production MCP server performs the following check before every tool
execution:

\begin{lstlisting}[language=TypeScript,caption={Production compose check (simplified)}]
function checkComposedPolicy(
  sandbox: SandboxPolicy,
  costGuard: CostGuard,
  cap: CapabilityClass,
): ExecutionResult {
  const sandboxOk = sandbox.allows(cap);
  const budgetOk  = !costGuard.isOverBudget();
  if (sandboxOk && budgetOk)  return 'allowed';
  if (!sandboxOk && budgetOk) return 'deniedSandbox';
  if (sandboxOk && !budgetOk) return 'deniedBudget';
  return 'deniedBoth';
}
\end{lstlisting}

This is the four-way match from the Lean \texttt{execute} definition
(Definition~\ref{def:composed}).  The composition theorem guarantees that
neither returning \texttt{allowed} when the sandbox denies nor returning
\texttt{allowed} when the budget is exceeded is possible.

\todo{Measurement: when the production test suite includes an integration test
exercising the full compose-check path (including the four-way discriminant),
cite the test here with commit hash.}

% ══════════════════════════════════════════════════════════════════════════════
\section{Axiom Audit}
\label{sec:axioms}

The \ATC{} mechanization introduces exactly three named axioms.  Every axiom
in a mechanized proof is a trust assumption; the axiom audit makes this
assumption explicit and documents its external discharge.

\paragraph{Audit gate.}
The Lean command \texttt{\#print axioms composition\_preserves\_both} lists
all axioms transitively reachable from the main theorem.  The output contains
exactly: \texttt{Classical.choice}, \texttt{Quot.sound},
\texttt{propext} (the standard Lean kernel axioms accepted universally), and
the three \ATC{} axioms listed below.  No additional opaque constants appear.

\paragraph{Axiom 1: \texttt{semiring\_merge\_commutative}.}
States that $\mergeadd_s$ is commutative for all five strategies $s \in
\mathcal{S}$.  \textbf{External validation}: Paper~3 (CRDT semiring,
ECOOP)~\cite{paper-3-crdt-ecoop2027} Theorem~1 verifies commutativity
empirically on 10,000/10,000 random pairs for all five strategies (commit
\texttt{66d58b58}).  Three of the five strategies also have closed-form
Lean proofs (\texttt{min\_plus\_commutative}, \texttt{max\_plus\_commutative},
\texttt{authority\_weighted\_commutative}; Theorems~23--25) that validate the
axiom concretely.

\paragraph{Axiom 2: \texttt{semiring\_merge\_associative}.}
States that $\mergeadd_s$ is associative for all five strategies.
\textbf{External validation}: Paper~3 Theorem~2; 10,000/10,000 pairs.
Three concrete Lean proofs validate the axiom for the
closed-form strategies (\texttt{min\_plus\_associative},
\texttt{max\_plus\_associative},
\texttt{authority\_weighted\_associative}; Theorems~26--28).

\paragraph{Axiom 3: \texttt{composition\_distributes\_over\_merge}.}
States that $\mergemul_s$ distributes over $\mergeadd_s$ (i.e.,
$(a \mergeadd_s b) \mergemul_s c = (a \mergemul_s c) \mergeadd_s (b \mergemul_s c)$)
for all strategies.
\textbf{External validation}: Paper~3 Theorem~3; 10,000/100,000 pairs.
This is the standard semiring distributivity law.  For the min-plus tropical
semiring it corresponds to the identity $\min(a,b)+c = \min(a+c,b+c)$, a
classical result~\cite{mohri2002semiring}.

\paragraph{Axiom design pattern.}
Stating semiring laws as axioms follows the same pattern as Paper~22's
\texttt{MSC} mechanization, where \texttt{solver\_functional} and
\texttt{cael\_causal\_well\_formed} capture properties of the production
runtime that the abstract model cannot internally witness.  In \ATC{}, the
axioms capture properties of the merge strategies that hold for all five
strategies but whose proofs for \texttt{domainOverride} and
\texttt{strictError} require case analysis on non-Nat carriers; stating
them as axioms---with concrete Lean proofs for the three closed-form
cases---is the appropriate mechanization discipline.

\todo{If future work closes the remaining two strategy axioms (domainOverride,
strictError) with closed-form Lean proofs, the axiom count drops from 3 to 3
kernel-only and this section can be updated accordingly.}

% ══════════════════════════════════════════════════════════════════════════════
\section{Related Work}
\label{sec:related}

\paragraph{Type-system composition.}
The foundational theory of type-system composition---most clearly expressed
in the work of Reynolds~\cite{wahbe1993efficient} on parametricity and of
Cardelli on subtyping lattices---establishes that well-typed composition
preserves semantic properties.
\todo{cite: Reynolds 1983 ``Types, Abstraction, and Parametric Polymorphism'', Cardelli--Wegner 1985}
The composition theorem in this paper instantiates this principle for a
specific two-component policy algebra.

\paragraph{Information-flow control.}
Sabelfeld and Myers~\cite{holoscript2026} survey the information-flow control
(IFC) literature, where the central challenge is composing security policies
across multiple principals without leakage.
\todo{cite: Sabelfeld--Myers 2003 ``Language-Based Information-Flow Security''}
IFC composition typically relies on non-interference, a property that our
sandbox preservation theorem closely parallels: a denied capability is the
analog of a secret that must not flow to a public output.  The budget
preservation theorem has no direct IFC counterpart---cost is a resource
bound, not an information-flow property---making the two-component composition
in \ATC{} a distinct contribution.

\paragraph{Verified sandbox systems.}
Wahbe et al.~\cite{wahbe1993efficient} introduced software fault isolation
(SFI) as an efficient sandboxing mechanism.  Yee et al.~\cite{yee2009native}
built Native Client, a verified SFI system for untrusted native code, with
formal safety arguments.  Haas et al.~\cite{haas2017wasm} extended these
ideas to WebAssembly.  Bosamiya et al.~\cite{cao2022wasmsandbox} demonstrated
provably-safe WebAssembly sandboxing at USENIX Security 2022.  These systems
prove sandbox properties individually; \ATC{} proves that combining a sandbox
with a cost guard does not weaken either.

\paragraph{Mechanized security proofs in Lean/Coq.}
\todo{cite: CompCert (Leroy 2009), CertiKOS (Gu et al. 2016), seL4 (Klein et al. 2009)}
The practice of mechanizing security properties in proof assistants has
produced significant results: CertiKOS~\cite{cao2022wasmsandbox} verified the
composition of OS-level security policies; seL4 provided the first
machine-checked proof of an OS kernel.  \ATC{} contributes a small but
complete composition result in a newer proof assistant (Lean~4) aimed
specifically at AI agent infrastructure.

\paragraph{Semiring frameworks for policy.}
Mohri~\cite{mohri2002semiring} established the theoretical foundations of
semiring-weighted automata, showing that the semiring laws (commutativity,
associativity, distributivity) are the load-bearing algebraic properties for
path-aggregation problems.  Green, Karvounarakis, and Tannen~\cite{green-semiring2007}
showed that semirings are the correct algebraic structure for database
provenance---a direct ancestor of HoloScript's provenance-merge strategies.
\ATC{}'s three semiring axioms are an instantiation of these foundations for
the five W.GOLD.189 merge strategies.

\paragraph{MCP and tool-use security.}
Hou et al.~\cite{hou2025mcp} survey the MCP ecosystem and identify capability
composition as an open security concern.  Greshake et al.~\cite{greshake2023}
and Zhan et al.~\cite{zhan2024} demonstrate injection and composition attacks
on tool-integrated LLM agents.  \ATC{} addresses the composition concern for
the specific case of sandbox+cost-guard composition, providing a mechanically
verified guarantee.

% ══════════════════════════════════════════════════════════════════════════════
\section{Conclusion}
\label{sec:conclusion}

We have presented \ATC{}, a Lean~4 mechanization proving that the composition
operator $\compose$ preserves both the sandbox-restriction invariant and the
budget-bounding invariant for AI agent tool execution.  The mechanization
comprises 23~theorems (0~\texttt{sorry}), 3~named axioms (all externally
validated by Paper~3's CRDT semiring suite), and a direct production mapping
to the HoloScript MCP server.

\paragraph{What is proved.}
\begin{itemize}
  \item No capability forbidden by the sandbox can be opened by composing with
        any cost guard (\texttt{sandbox\_preservation},
        \texttt{sandbox\_preservation\_for\_all\_cost\_guards}).
  \item No over-budget cost guard can be cleared by composing with any sandbox
        policy (\texttt{budget\_preservation},
        \texttt{budget\_preservation\_for\_all\_sandbox\_policies}).
  \item The composition operator forms a well-structured policy algebra:
        commutativity, identity, and annihilator laws all hold, enabling
        equational reasoning about policy chains.
  \item The four-valued execution result faithfully reflects both invariants in
        the observable outcome (\texttt{execute\_preserves\_sandbox},
        \texttt{execute\_preserves\_budget},
        \texttt{allowed\_implies\_both\_permitted}).
\end{itemize}

\paragraph{Open: circuit-level enforcement layer.}
The composition theorem covers the logical structure of policy composition.
Layer~2 (CAEL hash-chain audit of every composed execution) and Layer~3
(simulation-contract oracle) remain open for Lean mechanization and are
targeted by Paper~4~\cite{paper-4-sandbox-usenix} and the companion
capstone~\cite{trust-by-construction}.

\paragraph{Open: closing the remaining strategy axioms.}
The two strategies without closed-form Lean proofs
(\texttt{domainOverride}, \texttt{strictError}) could be given concrete
mechanizations with a richer carrier type; this would reduce the axiom count
from three to zero.  We leave this for future work.

\paragraph{Open: multi-component composition.}
The current theorem handles exactly two components (sandbox $\compose$
cost-guard).  An $n$-component composition theorem, admitting arbitrarily
many policy layers stacked via $\compose$, would generalize the result for
multi-tenant fleet scenarios and is a natural extension.

% ══════════════════════════════════════════════════════════════════════════════
\begin{acks}
This work is part of the HoloScript simulation-as-proof research program.
The Lean mechanization builds on the provenance-semiring proofs in Paper~3
(CRDT).  The composition operator design was informed by Paper~4 (Sandbox)
and Paper~8 (Unified Transform).
\end{acks}

% ══════════════════════════════════════════════════════════════════════════════
\bibliography{holoscript}
\bibliographystyle{ACM-Reference-Format}

\end{document}
